\begin{textsample}{POS Dim 4 – persona\_gemini – Score 14.00 – t421\_gemini.txt}  \label{ex:f4_pos_049}
In Baja California Sur, Mexico, the coral \textbf{reef} of Cabo Pulmo is home to a diverse array of marine \textbf{life}, including \textbf{fish}, crustaceans, mollusks, \textbf{birds}, and \textbf{mammals}, many of which are at risk.
Twenty-five years ago, this \textbf{area} was threatened by overfishing and irresponsible tourism.
In response, local communities organized, an effort that led to Cabo Pulmo ’s declaration as a National Marine Park in 1995.
The results were remarkable. \textbf{Fish} populations increased by over 400 %, and migratory \textbf{species} such as \textbf{whale} \textbf{sharks}, mantas, humpback \textbf{whales}, \textbf{turtles}, and \textbf{sharks} returned to the \textbf{area}.
The local communities successfully shifted their economies to ecotourism.
However, in 2008, a new \textbf{threat} emerged.
A project by the company Hansa Urbana, known as"Cabo Cortés,"posed a danger of pollution and \textbf{habitat} destruction to the \textbf{reef}.
A new movement rose to defend the park, and campaigns with Greenpeace, which included a petition signed by 220,000 people, led to the project 's cancellation in 2012.
The story of Cabo Pulmo highlights the impact of people power.
It is an example of why we must work to protect 30 % of our \textbf{oceans} through the creation of \textbf{ocean} \textbf{sanctuaries} and the implementation of a Global Ocean Treaty.

% matched lemmas: area, bird, fish, habitat, life, mammal, ocean, reef, sanctuary, shark, specie, threat, turtle, whale
\end{textsample}
